\Card{11-01}{Commit to the Moon}
{Planned TLI ignition 002:44:26; 320 s burn. [F1, Table 1-1]}
{Armstrong commands, Aldrin reads the burn pad, and Collins checks the
decimal place. The Saturn's third stage will send Columbia and Eagle
toward the Moon. Do not mistake feet per second for metres per second.}
{Convert the planned \(10{,}451.2\) ft/s velocity change to m/s.
Find its average acceleration over 5 min 20 s. Compute the cutoff GET.}
{Convert minutes by hand. On C/D multiply velocity by 0.3048, then divide
by elapsed seconds: \(\bar a=\Delta v/t\). Add the duration to GET on paper.
This is a whole-burn average, not constant actual thrust acceleration.}
{Write m/s, m/s\(^{2}\), and h:m:s on the pad. Missing or incompatible
units are RED. The modeled burn is commanded in the forward direction.}
{AMBER: guidance flags a targeting discrepancy; charge 1 reserve for a
corrective update. RED: the departure targeting test fails; charge 2 for
a simulator abort-and-replan. A depleted reserve here means a safe
Earth-orbit abort, not a lunar landing.}
{\(\Delta v=10{,}451.2(0.3048)=3185.52576\) m/s, so report
\(\boxed{3.19\times10^3\text{ m/s}}\).
\(t=5(60)+20=320\) s;
\(\bar a=3185.52576/320=\boxed{9.95\text{ m/s}^2}\).
Cutoff \(002{:}44{:}26+00{:}05{:}20=\boxed{002{:}49{:}46}\).
An estimate near \(10^4\) ft/s and \(3\times10^3\) m/s catches a factor-ten error.}

\Card{11-02}{Clear the spent stage}
{Planned evasive maneuver 004:39:44.9; 2.8 s. [F1, Table 1-1]}
{Collins separates the docked spacecraft from the spent stage.
The maneuver pad has two nonzero vector components, not two separate
burns. Mission Control wants a magnitude before authorizing ignition.}
{For \((\Delta v_x,\Delta v_y,\Delta v_z)=(5.1,0,19.0)\) ft/s,
find total velocity change in ft/s and m/s. A training model supplies
average acceleration \(2.14\) m/s\(^{2}\); find burn duration.}
{Square on A (or multiply C/D), add the squares with pencil, take the
square root on A/D:
\(\Delta v=\sqrt{\Delta v_x^2+\Delta v_y^2+\Delta v_z^2}\).
Convert on C/D and divide by \(\bar a\).}
{Use the vector magnitude, not \(5.1+19.0\).
Preserve the supplied component directions on the command pad.
The 2.14 m/s\(^{2}\) acceleration is invented training telemetry.}
{AMBER: separation is slower or faster than expected; charge 1 for a
tracking update. RED: the safe-separation prediction fails; charge 2
for an emergency avoidance rehearsal before continuing.}
{\(5.1^2+19^2=26.01+361=387.01\).
\(\Delta v=\sqrt{387.01}=19.67257\) ft/s
\(\rightarrow\boxed{19.7\text{ ft/s}}\).
\(\Delta v=19.67257(0.3048)=\boxed{6.00\text{ m/s}}\).
\(t=5.996199/2.14=\boxed{2.80\text{ s}}\), matching the
planned 2.8 s to slide-rule precision.
Use the table's vector / required-\(\Delta v\) data here, not the
different 15.6 ft/s value in its \(\Delta V_C\) column.}

\Card{11-03}{Capture lunar orbit}
{Planned LOI1 ignition 075:54:28.4; 358.9 s. [F1, Table 1-1]}
{Columbia passes behind the Moon. Houston cannot coach this burn in real
time. Collins reads a retrograde pad; the crew's arithmetic determines
whether the simulator accepts lunar capture.}
{Use components \((-2891.8,-433.1,+20.4)\) ft/s to find
\(\Delta v\) in ft/s and m/s. Find average acceleration for
5 min 58.9 s. Compute cutoff GET.}
{Use A/D for squares and roots, pencil for addition, C/D for unit
conversion and division. Use \(\bar a=\Delta v/t\).
Clock arithmetic is sexagesimal; keep the tenths of a second.}
{Keep all three component signs. A correct magnitude with reversed
retrograde direction is RED. This scalar exercise does not compute
an orbit or prove capture in the real world.}
{AMBER: tracking predicts an off-target lunar orbit; charge 1 for a
trim update. RED: the simulator rejects capture; charge 2 for a
contingency return rehearsal. Reserve exhaustion means abort the
landing and use the game's safe-return ending.}
{\(\Sigma \Delta v_i^2=8{,}550{,}499.01\) (ft/s)\(^{2}\).
Square root \(=2924.12363\) ft/s
\(\rightarrow\boxed{2.92\times10^3\text{ ft/s}}\).
Multiply by 0.3048 to get \(891.273\) m/s
\(\rightarrow\boxed{891\text{ m/s}}\).
\(t=358.9\) s and \(\bar a=891.273/358.9
=\boxed{2.48\text{ m/s}^2}\).
Cutoff \(075{:}54{:}28.4+00{:}05{:}58.9
=\boxed{076{:}00{:}27.3}\).
The table rounds required \(\Delta v\) to 2924.1 ft/s;
its separate \(\Delta V_C=2914.8\) ft/s is not substituted.}

\Card{11-04}{Work the lunar clock}
{Lunar-orbit operations between LOI2 and LM separation. [F1, p. 1-3]}
{Houston schedules a daylight observation and a radio call.
Collins must distinguish orbital night from loss of radio coverage:
they are different scheduling constraints.}
{An average orbit lasts 118.2 min, including 47 min of night.
Find daylight minutes, average angular rate in degrees/min,
and the percentage of an orbit with MSFN coverage if coverage is
72 min. Is a 65 min daylight task feasible in one daylight interval?}
{Subtract with pencil, then use C/D for \(360/118.2\) and
\(100(72/118.2)\). Compare the task with daylight duration.
Use the supplied full period, not the plan's rounded \(3^\circ/\)min.}
{The 65 min task must fit daylight. Do not treat 72 min of coverage as
72 min of daylight, or claim that their overlap is known.
The task duration is a training scenario.}
{AMBER: a mistimed observation loses useful coverage; charge 1 for
rescheduling. RED: the crew confuses the two windows and misses the
handover; charge 2 for a complete scheduling rebuild.}
{Daylight \(=118.2-47=\boxed{71.2\text{ min}}\).
Angular rate \(=360/118.2=\boxed{3.05^\circ/\text{min}}\).
Coverage \(=100(72/118.2)=\boxed{60.9\%}\).
\(\boxed{\text{GO}}\) for the daylight task:
\(71.2-65=\boxed{6.20\text{ min}}\) spare.
The plan gives approximately \(3^\circ/\)min; using the more precise
period explains the small difference. There is insufficient information
to require continuous radio contact throughout the task.}

\Card{11-05}{Lower Eagle's perilune}
{Planned DOI 101:38:48, 28.5 s; PDI 102:35:13. [F1, Table 1-2]}
{Aldrin checks descent-orbit insertion while Collins tracks Eagle.
Armstrong wants to know how much coast time remains after engine cutoff
before powered descent begins.}
{Use the table's required \(\Delta v=70.0\) ft/s.
Convert it to m/s; compute average acceleration over 28.5 s.
Find DOI cutoff GET and the coast interval from cutoff to PDI.}
{C/D handles conversion and \(\Delta v/t\).
Add the burn on paper, then subtract cutoff from PDI, borrowing
60 seconds or minutes when necessary.}
{Use required \(\Delta v\), not a reconstruction from the printed
components: this table's DOI components do not exactly reproduce 70.
The game does not resolve that source discrepancy.}
{AMBER: the coast briefing is late; charge 1 for an expedited pad
review. RED: a descent-readiness check is missed; charge 2 for a
full rehearsal. If reserve is exhausted, cancel PDI and rendezvous
with Columbia in the game's safe-return ending.}
{\(70(0.3048)=\boxed{21.3\text{ m/s}}\).
\(\bar a=21.336/28.5=\boxed{0.749\text{ m/s}^2}\).
Cutoff \(101{:}38{:}48+00{:}00{:}28.5
=\boxed{101{:}39{:}16.5}\).
PDI minus cutoff:
\(102{:}35{:}13-101{:}39{:}16.5
=\boxed{00{:}55{:}56.5}\).
The printed components \(67.46,-28.68,-12.51\) ft/s would give about
74.4 ft/s, so the specified 70.0 ft/s input is deliberate.}

\Card{11-06}{Land or wave off}
{Planned PDI 102:35:13 to touchdown 102:47:11. [F1, pp. 1-9, 3-69]}
{Armstrong sees a rough patch ahead. Aldrin calls altitude and fuel.
Choose the longest safe diversion before the checkpoint commits the
landing. This last-minute hover model is fictional, not Apollo telemetry.}
{First find touchdown GET and average whole-descent acceleration from
\(\Delta v=6766\) ft/s and 11 min 58 s.
At the terminal gate, altitude is 120 m and downward speed is held at
2.00 m/s. Fuel is 60.0 kg and consumption is 0.700 kg/s.
Find descent time, fuel endurance, and maximum extra hover time.
Choose the longest safe diversion from 10, 20, or 30 s; find fuel left.}
{Use C/D to convert velocity, divide for average acceleration,
divide altitude by speed, and divide fuel by flow.
Subtract endurance minus descent time. For a diversion \(d\),
\(m_{\rm left}=60-0.700(120/2+d)\).}
{Choose the longest listed diversion that leaves nonnegative fuel.
The terminal constant-speed model is separate from the whole-descent
average; do not use that average as net terminal acceleration.}
{AMBER: the landing pad needs a last-second correction; charge 1 and
land at a less desirable site. RED: fuel prediction rejects the
landing; charge 2 and rehearse an immediate wave-off. If reserve is
exhausted here, no landing occurs; end with a safe lunar-orbit abort.}
{Touchdown \(=\boxed{102{:}47{:}11}\).
\(t=718\) s; \(\Delta v=6766(0.3048)=2062.2768\) m/s;
\(\bar a=\boxed{2.87\text{ m/s}^2}\).
Terminal descent \(=120/2=\boxed{60.0\text{ s}}\).
Endurance \(=60/0.700=\boxed{85.7\text{ s}}\);
extra hover \(=85.7143-60=\boxed{25.7\text{ s}}\).
Choose \(\boxed{20\text{ s}}\).
Fuel left \(=60-0.700(80)=\boxed{4.00\text{ kg}}\).
A 30 s diversion would require 63 kg and is NO-GO.
Ten seconds is safe but not the requested longest safe choice.}

\Card{11-07}{A walk with a turnaround time}
{Planned EVA duration 2 h 40 min. [F1, p. 1-2]}
{Armstrong and Aldrin share the science work, but each carries an
independent oxygen budget. The flight director will not spend the
walk-back reserve on one more rock.}
{For each walker, oxygen starts at 1.80 kg; modeled use is
0.480 kg/h and protected reserve is 0.300 kg.
Find use and remaining oxygen over the 160 min EVA.
Find maximum allowed EVA duration and schedule margin.
Egress and return consume 28 min of the planned EVA; each science
station takes 12 min. How many complete stations fit?}
{C/D: convert minutes to hours, multiply by flow,
and divide usable oxygen by flow.
Pencil: subtract reserves and non-science time.
Round station count down, never up.}
{Both walkers must retain at least 0.300 kg.
The protected reserve is not available for science. These suit quantities
and task times are invented for play.}
{AMBER: Houston calls an early turnaround; charge 1 and record
``reduced science.'' RED: the oxygen or station plan violates a return
margin; charge 2 and record ``EVA curtailed.'' Correct the budget
before ascent preparations.}
{EVA \(=160/60=2.66667\) h.
Use \(=0.480(2.66667)=\boxed{1.28\text{ kg per walker}}\).
Remaining \(=1.80-1.28=\boxed{0.520\text{ kg}}\), so GO.
Allowed duration \(=(1.80-0.300)/0.480=3.125\) h
\(=\boxed{187.5\text{ min}}\) (3 h 7 min 30 s).
Margin \(=187.5-160=\boxed{27.5\text{ min}}\).
Stations \(=\lfloor(160-28)/12\rfloor=\boxed{11}\).
Do not divide one walker's tank by the combined two-person use rate.}

\Card{11-08}{Rise, meet, and dock}
{Planned ascent 124:23:26, 438 s; CSI 125:21:19.1. [F1, Table 1-2]}
{Eagle leaves the surface. Collins will bring Columbia to the rendezvous
only after Aldrin reads back ascent timing and a separate close-range
braking command.}
{For ascent \(\Delta v=6060\) ft/s, find m/s and average acceleration
over 7 min 18 s. Find insertion GET and time until CSI.
At the separate docking exercise, range is 40.0 m and closure speed
0.200 m/s. Find time to contact without braking and a burn duration
to cancel closure using \(0.0500\) m/s\(^{2}\).}
{Use C/D for conversion and division, pencil for GET.
For docking, \(t_{\rm contact}=r/v\) and \(t_{\rm brake}=\Delta v/a\).
The nominal burn is initiated immediately; stop the modeled closure
before contact and then let the referee narrate docking.}
{Command braking opposite to closure. Wrong direction is RED.
The close-range model is not the CSI maneuver and does not model
the actual rendezvous orbit.}
{AMBER: Collins delays docking while the crew corrects its pad; charge 1.
RED: the docking safety predictor flags a collision; charge 2 for an
emergency separation rehearsal. Exhaustion after landing means the
simulation ends with rendezvous unresolved, not an invented safe rescue.}
{\(\Delta v=6060(0.3048)=\boxed{1.85\times10^3\text{ m/s}}\).
\(\bar a=1847.088/438=\boxed{4.22\text{ m/s}^2}\).
Insertion \(124{:}23{:}26+00{:}07{:}18
=\boxed{124{:}30{:}44}\).
CSI coast \(125{:}21{:}19.1-124{:}30{:}44
=\boxed{00{:}50{:}35.1}\).
Time to contact \(=40/0.200=\boxed{200\text{ s}}\).
Braking \(=0.200/0.0500=\boxed{4.00\text{ s}}\), opposite closure.
During immediate uniform braking, travel is only
\(\tfrac12(0.200)(4)=0.400\) m, safely below the 40 m range.}

\Card{11-09}{Bring Columbia home}
{Planned TEI 135:24:33.8, 149.4 s; EI 195:03:27. [F1]}
{After docking and LM jettison, Collins checks the Earth-return pad.
Houston supplies signed components and the planned entry-interface
time, not a promise that every computed trajectory reaches Earth.}
{For components \((3213.3,705.0,-138.8)\) ft/s, find vector magnitude
in ft/s and m/s, and average acceleration over 2 min 29.4 s.
Find cutoff GET and coast time from cutoff to entry interface.}
{Use A/D to square and take a root; pencil to add;
C/D to convert and divide. Convert duration before adding GET.
Subtract the cutoff time, not ignition, from entry interface.}
{Retain the signed components on the burn command. An incorrect
direction is RED even with the right magnitude. Coast time is a
schedule calculation, not an orbital prediction.}
{AMBER: Earth-return targeting needs a trim; charge 1.
RED: the entry predictor rejects the pad; charge 2 for a return-burn
replan. Exhaustion means an unresolved return simulation.}
{\(\Sigma \Delta v_i^2=10{,}841{,}587.33\) (ft/s)\(^{2}\).
\(\Delta v=3292.65658\) ft/s
\(\rightarrow\boxed{3.29\times10^3\text{ ft/s}}\).
Multiply by 0.3048: \(\boxed{1.00\times10^3\text{ m/s}}\)
(guard value 1003.60172).
\(\bar a=1003.60172/149.4=\boxed{6.72\text{ m/s}^2}\).
Cutoff \(=\boxed{135{:}27{:}03.2}\).
Coast \(195{:}03{:}27-135{:}27{:}03.2
=\boxed{59{:}36{:}23.8}\).
The plan's rounded required magnitude is 3292.7 ft/s.}

\Card{11-10}{Entry corridor and recovery}
{Planned EI 195:03:27; detailed splashdown 13:55 later. [F1, p. 3-135]}
{The crew checks entry geometry before separating for the final descent.
Houston's training predictor supplies a downward vertical speed and a
horizontal speed. It is a diagnostic exercise, not an atmospheric model.}
{At the training gate, horizontal speed is 11.0 km/s and downward
vertical speed is 1.25 km/s. Find flight-path angle
\(\gamma=-\arctan(v_{\rm down}/v_{\rm horizontal})\), speed magnitude,
and kinetic energy per unit mass in MJ/kg.
Check the modeled corridor \(-6.8^\circ\leq\gamma\leq-6.2^\circ\)
and energy limit 63.0 MJ/kg. The detailed plan gives splashdown at
EI plus 13 min 55 s. Find splashdown GET and confirm the interval.}
{C/D gives the speed ratio; interpolate the tangent table on the
reference sheet or use T. Use A/D for speed.
\(e=v^2/2\): a squared speed in (km/s)\(^{2}\) gives MJ/kg directly.
Use pencil for clock subtraction.}
{Both corridor and energy checks must be GO. Positive \(\gamma\)
is an upward trajectory and is RED. These limits are game thresholds,
not authenticated Apollo flight rules.}
{AMBER: the predictor requests a corrected pad; charge 1.
RED: a shallow angle predicts skip-out, a steep angle predicts
excessive deceleration, or an energy error flags heating uncertainty;
charge 2. Do not enter until corrected. With reserve intact, narrate
parachutes and recovery; otherwise the return simulation is unresolved.}
{Ratio \(=1.25/11=0.113636\).
Table interpolation gives about \(6.48^\circ\);
\(\boxed{\gamma=-6.48^\circ}\).
\(v=\sqrt{11^2+1.25^2}=\boxed{11.1\text{ km/s}}\).
\(e=(121+1.5625)/2=\boxed{61.3\text{ MJ/kg}}\).
\(\boxed{\text{GO on both gates}}\).
Splashdown \(195{:}03{:}27+00{:}13{:}55
=\boxed{195{:}17{:}22}\).
Subtracting EI confirms \(\boxed{00{:}13{:}55}\).
The summary rounds this to 195:17; the game uses the more precise
detailed timeline rather than inventing zero seconds.}

\ifanswers\else
\clearpage
\section*{Apollo 13: change the mission, save the crew}
\begin{panel}{Advanced mission / fresh command reserve: 4}
Crew: Lovell, Haise, Swigert. Narrate successful launch, docking, and
outbound coast. At GET 055:54:53 the service-module oxygen-tank accident
occurs. The Moon is now a turning point, not a destination.
Your victory condition is recovery of all three crew after ten
corrected checkpoints, with command reserve still nonnegative.
\end{panel}

Unlike the Apollo 11 campaign, this mission requires independent
full-horizon consumable audits, a logarithmic rocket-equation
calculation, cartridge-count rounding, signed manual-navigation geometry,
and a linear entry-correction model.

\begin{tabularx}{\linewidth}{@{}lX@{}}
\toprule
Historical anchor & Event used in play\\
\midrule
055:54:53 & Oxygen-tank rupture; activate the LM lifeboat.\\
Near 061:30 & LM propulsion restores the desired free-return path.\\
079:27:39 & PC+2 begins after the lunar swingby.\\
Near 105:19 & Later manual correction; about a 15 s burn.\\
Final correction & A 22.4 s burn motivates the entry-predictor exercise.\\
141:30:02 & Aquarius jettison; Odyssey must already be ready.\\
142:54:41 & Historical splashdown and crew recovery.\\
\bottomrule
\end{tabularx}

\textbf{Planning epoch.} Consumable cards 13-03, 13-04, and 13-06 all
start their accounting at 058:00:00, regardless of when they are discussed
in the narrative. Plan the entire remaining horizon only once per pool.
Copy the time horizon to scrap paper and reuse it.

\textbf{Provenance.} The supplied PDF and text are both Apollo 11
resources. NASA's Apollo 13 catalog supplies these rescue event anchors.
All game inventories, simplified burn models, and failure thresholds
are invented and marked as training data.

\textbf{Before committing.} The commander checks direction, the LM
pilot checks the protected reserve, and the CM pilot checks time and
units. Any player may catch an error for free before ``commit.''
Start card 13-01 with a fresh ledger; do not carry Apollo 11 resources.
\fi

\Card{13-01}{The alarm: power the lifeboat}
{Historical tank rupture 055:54:53. [N1]; this activation model is invented.}
{Lovell takes command, Swigert powers down Odyssey, and Haise prepares
Aquarius. The explosion occurs regardless of your previous success.
Your objective changes from landing to bringing all three crew home.}
{In a separate emergency CM battery pool, usable energy is 2.40 kWh.
Load is 1.20 kW and LM activation takes 22 min.
Find initial endurance, activation energy, remaining energy,
and endurance remaining at that load. Find activation completion GET.
Is activation possible before this pool is exhausted?}
{Use C/D: \(t=E/P\), \(E=Pt\), minutes divided by 60.
Subtract energy on paper; add 22 min to the alarm time.
This emergency pool is not the protected entry battery on card 13-10.}
{Activate Aquarius immediately and transfer the crew.
Keeping Odyssey at this load indefinitely is NO-GO.
The explosion is a scripted historical event, never a penalty for bad math.}
{AMBER: transition checklist delays threaten battery margin; charge 1.
RED: the power-down budget is rejected; charge 2 for a contingency
transfer rehearsal. Reserve exhaustion means the rescue simulation fails.}
{Initial endurance \(=2.40/1.20=\boxed{2.00\text{ h}}\).
Activation \(=22/60=0.366667\) h;
energy \(=1.20(0.366667)=\boxed{0.440\text{ kWh}}\).
Remaining \(=2.40-0.440=\boxed{1.96\text{ kWh}}\).
Remaining endurance \(=1.96/1.20=\boxed{1.63\text{ h}}\)
(98 min).
Completion \(055{:}54{:}53+00{:}22{:}00
=\boxed{056{:}16{:}53}\).
\(\boxed{\text{GO}}\): 22 min is less than 120 min.}

\Card{13-02}{Restore a free return}
{Historical LM return maneuver near GET 061:30; magnitude 11.6 m/s. [N1]}
{The docked stack is not on the desired return path.
Haise reads a two-axis training pad while Lovell checks its direction.
This burn is distinct from the later PC+2 maneuver.}
{Training components are \((+11.0,+3.70)\) m/s.
Find total \(\Delta v\) and burn time using average acceleration
0.400 m/s\(^{2}\).
For a modeled constant thrust 12,000 N, \(I_{\rm sp}=305\) s,
and \(g_0=9.81\) m/s\(^{2}\), find mass flow and propellant used.}
{A/D handles vector magnitude. C/D:
\(t=\Delta v/\bar a\),
\(\dot m=F/(I_{\rm sp}g_0)\), \(m_{\rm used}=\dot m t\).
Keep intermediate readings on the sheet; do not repeatedly round them.
The fixed average acceleration is a supplied approximation.}
{Preserve both positive components. For this isolated maneuver,
usable propellant is 150 kg and at least 20 kg must remain.
Check the fuel gate; it is not the PC+2 tank inventory.}
{AMBER: return tracking predicts a correction; charge 1.
RED: a wrong-way or underbudgeted burn fails the free-return
predictor; charge 2 for a replanning rehearsal. No lunar landing is
available as an alternative ending.}
{\(\Delta v=\sqrt{11^2+3.7^2}
=\sqrt{134.69}=\boxed{11.6\text{ m/s}}\).
\(t=11.605602/0.400=\boxed{29.0\text{ s}}\).
\(\dot m=12{,}000/(305\cdot9.81)
=\boxed{4.01\text{ kg/s}}\).
\(m_{\rm used}=4.010629(29.014005)
=\boxed{116\text{ kg}}\) (guard value 116.364).
Remaining \(=150-116.364=\boxed{33.6\text{ kg}}\).
\(\boxed{\text{GO}}\), above the 20 kg reserve.
Mass and thrust values are training inputs, not the historical burn pad.}

\Card{13-03}{Three people, one oxygen budget}
{Game accounting epoch 058:00:00; historical splashdown 142:54:41. [N1]}
{Aquarius was designed for a different mission. Haise must choose a
conservation profile for the entire remaining flight. The same exact
time horizon will be used on the electrical and water cards.}
{Compute remaining flight time \(T\) in h:m:s and decimal hours.
LM usable oxygen is 12.0 kg with a protected 2.00 kg reserve.
Compare ``conserve'' use 0.0350 kg/(person h) against ``normal''
0.0500 kg/(person h), for three people.
Calculate total use and remaining oxygen for each; choose a safe profile.}
{Subtract GET on paper.
On C/D calculate \(T=84+54/60+41/3600\) and
\(m=3rT\). Subtract each use from 12.0 kg.
Keep \(T\) to at least four digits on paper for later cards.}
{Remaining oxygen must be at least 2.00 kg.
Do not confuse per-person use with total cabin use.
Carry the chosen profile to the mission ledger; all rates are
fictional planning averages.}
{AMBER: an incorrect budget prompts tighter rationing; charge 1.
RED: the selected profile cannot sustain the protected oxygen margin;
charge 2 for an emergency conservation replan.}
{\(142{:}54{:}41-058{:}00{:}00
=\boxed{84{:}54{:}41}\);
\(T=\boxed{84.9113889\text{ h}}\), normally read as 84.9 h.
Conserve: \(3(0.0350)T=\boxed{8.92\text{ kg}}\);
remaining \(\boxed{3.08\text{ kg}}\) (3.084304), GO.
Normal: \(3(0.0500)T=\boxed{12.7\text{ kg}}\);
remaining \(\boxed{-0.737\text{ kg}}\) (minus 0.736708), NO-GO.
Choose \(\boxed{\text{conserve}}\).
Its margin above the protected 2 kg is 1.084304 kg.}

\Card{13-04}{A dark, cold electrical plan}
{Same 058:00:00--142:54:41 horizon as 13-03.}
{Swigert asks which loads can remain off.
Haise proposes a low-power Aquarius profile; Lovell insists on
keeping entry energy protected rather than spending it on comfort.}
{Use \(T\) from 13-03. LM energy pool is 18.0 kWh; retain 2.50 kWh.
Non-idle operations cost 0.600 kWh in addition to idle use.
Compare idle loads 0.160 kW (low) and 0.240 kW (normal).
Find total use and remaining energy for each, and maximum permissible
idle power. Choose a profile.}
{C/D: \(E_{\rm use}=PT+0.600\).
Pencil: \(E_{\rm left}=18-E_{\rm use}\).
Then \(P_{\max}=(18-2.50-0.600)/T\).
Energy is kWh; power is kW. Do not multiply power by seconds here.}
{The low-power plan must retain at least 2.50 kWh.
The 0.600 kWh operations allowance already includes modeled LM burns;
do not charge them again. This pool is separate from CM entry energy.}
{AMBER: incorrect switching instructions force load shedding; charge 1.
RED: the chosen profile predicts depletion before recovery; charge 2
and rewrite the electrical plan. Cold and dim conditions are narrative
effects, not additional hidden resource deductions.}
{Low use \(=0.160(84.9113889)+0.600
=\boxed{14.2\text{ kWh}}\) (14.185822);
remaining \(\boxed{3.81\text{ kWh}}\).
Normal use \(=0.240T+0.600=\boxed{21.0\text{ kWh}}\);
remaining \(\boxed{-2.98\text{ kWh}}\).
\(P_{\max}=14.9/84.9113889
=\boxed{0.175\text{ kW}}\) (175.477 W).
Choose \(\boxed{\text{low}}\).
At 0.160 kW the reserve margin is 1.314178 kWh above 2.50.}

\Card{13-05}{PC+2: shorten the trip}
{Historical PC+2 ignition 079:27:39; burn 263.4 s, \(\Delta v=262\) m/s. [N1]}
{After passing the Moon, Aquarius must shorten the return.
This is the advanced slide-rule challenge: solve the rocket equation
backward, then compare the predicted duration with the historical anchor.}
{For a fictional mass model, \(m_0=30{,}000\) kg,
effective exhaust velocity \(v_e=3000\) m/s, target \(\Delta v=262\) m/s,
and average propellant flow 9.525 kg/s.
Find final mass, propellant used, and burn duration.
Find historical scheduled cutoff using 263.4 s, not your rounded prediction.}
{Use C/D and L:
\(x=\Delta v/(2.302585v_e)\),
\(m_f=m_0\,10^{-x}\),
\(m_{\rm used}=m_0-m_f\), \(t=m_{\rm used}/\dot m\).
On L, find the mantissa \(1-x\), read \(10^{1-x}\) on D,
then divide by 10 to obtain \(10^{-x}\).
Subtraction uses pencil; retain extra readable digits before subtracting.}
{The isolated burn's usable propellant is 2700 kg; retain 150 kg.
Mass must decrease, so \(m_f<m_0\). A wrong exponent sign is RED.
The historical time is an anchor; the masses and flow are invented.}
{AMBER: duration disagrees with the acceleration plan; charge 1 for a
throttle-pad correction. RED: the burn budget or exponent is wrong;
charge 2 for an emergency targeting rehearsal.}
{\(x=262/(2.302585\cdot3000)=0.03792838\).
\(10^{-x}=0.91637159\).
\(m_f=30{,}000(0.91637159)
=\boxed{27{,}491.15\text{ kg}}\) (slide-rule reading about 27,490).
Used \(=30{,}000-27{,}491.15
=\boxed{2508.85\text{ kg}}\), report \(2.51\times10^3\) kg.
\(t=2508.85237/9.525=\boxed{263\text{ s}}\)
(263.39657 s).
Remaining burn pool \(=2700-2508.85
=\boxed{191\text{ kg}}\), GO.
Historical cutoff \(079{:}27{:}39+00{:}04{:}23.4
=\boxed{079{:}32{:}02.4}\).
Allow the normal 1\% band on propellant despite subtractive sensitivity.}

\Card{13-06}{Cooling water versus drinking water}
{Same full return horizon \(T=84.9113889\) h; planning epoch 058:00:00.}
{Haise balances cooling against the crew's drinking allowance.
Choose a modeled cooling profile without borrowing the water
protected for entry and recovery.}
{Water starts at 100 kg; protect 12.0 kg.
Each of three crew uses 0.0600 kg/h for drinking.
Compare cooling 0.900 kg/h (normal) with 0.700 kg/h (reduced).
Find combined rates, total use, and water left for both profiles.
Find the maximum allowable cooling rate, then choose a safe profile.}
{C/D: \(r_{\rm total}=r_{\rm cool}+3r_{\rm drink}\),
\(m_{\rm use}=r_{\rm total}T\).
Subtract from 100 on paper.
\(r_{\rm cool,max}=(100-12)/T-3r_{\rm drink}\).}
{Keep at least 12.0 kg. Both proposed cooling settings are assumed
thermally permissible in this game; no thermal engineering calculation
is implied. The drinking allowance is not optional.}
{AMBER: a bad ration sheet prompts an uncomfortable reduced-cooling
replan; charge 1. RED: the selected profile consumes protected water;
charge 2 for emergency ration planning. Do not invent additional
water losses for narrative delays.}
{Drinking rate \(=3(0.0600)=0.180\) kg/h.
Normal total \(=\boxed{1.08\text{ kg/h}}\);
use \(=1.08T=\boxed{91.7\text{ kg}}\);
left \(\boxed{8.30\text{ kg}}\), NO-GO.
Reduced total \(=\boxed{0.880\text{ kg/h}}\);
use \(=0.880T=\boxed{74.7\text{ kg}}\);
left \(\boxed{25.3\text{ kg}}\), GO.
Maximum cooling \(=88/T-0.180
=\boxed{0.856\text{ kg/h}}\).
Choose \(\boxed{\text{reduced}}\).
Its margin above reserve is 13.277978 kg.}

\Card{13-07}{Make the square cartridge fit}
{Training checkpoint GET 090:00:00; CO\(_2\) planning to 134:00:00.}
{The astronauts need a cartridge adapter. The referee narrates its
construction from available materials; the mathematics determines
how many compatible cartridges and how soon they are needed.
This is not a real breathing-system construction guide.}
{Three people produce 0.900 kg of CO\(_2\) per person per day.
Existing compatible absorbent has 0.600 kg capacity remaining.
Each adapted spare removes 1.80 kg.
Find hourly production, the load for the next 44 h, and the minimum
integer number of spares needed in addition to existing capacity.
Find capacity margin, time until existing capacity runs out,
and the adapter deadline GET. Assembly takes 2.00 h: find time slack.}
{Use C/D for rate conversion and load.
Compute \(n=\lceil(m_{\rm load}-0.600)/1.80\rceil\): round UP.
Deadline \(t=0.600/r_{\rm hourly}\).
Add deadline to 090:00:00 and subtract assembly time on paper.}
{Finish the adapter before existing capacity expires and allocate enough
spares. Timing this card uses decimal endurance bands and clock bands;
cartridge count must be the exact minimum integer.}
{AMBER: a miscount requires inventory review; charge 1.
RED: the predicted scrubber gap would allow CO\(_2\) accumulation;
charge 2 for a rushed adapter rehearsal. The referee must not allow
the game to proceed with a knowingly inadequate scrubber plan.}
{Production \(=3(0.900)/24=\boxed{0.1125\text{ kg/h}}\).
Load \(=0.1125(44)=\boxed{4.95\text{ kg}}\).
\(n=\lceil(4.95-0.600)/1.80\rceil
=\lceil2.41667\rceil=\boxed{3\text{ spares}}\).
Capacity margin \(=0.600+3(1.80)-4.95
=\boxed{1.05\text{ kg}}\).
Existing endurance \(=0.600/0.1125
=\boxed{5.33333\text{ h}}\) (5 h 20 min).
Deadline \(=\boxed{095{:}20{:}00}\).
Assembly slack \(=5.33333-2=\boxed{3.33333\text{ h}}\)
(3 h 20 min). GO; finish by 092:00:00 if starting now.}

\Card{13-08}{Navigate without the usual computer}
{Historical correction near GET 105:19 used about a 15 s burn. [N1]}
{Lovell uses the visible Earth as an attitude reference.
Haise works the training pad by hand. There is no die roll for aim:
your signed vector and angle decide whether the predictor accepts it.}
{A fictional two-axis correction requires
\((\Delta v_x,\Delta v_y)=(-2.30,+0.400)\) m/s.
Find magnitude and burn time at average acceleration 0.155 m/s\(^{2}\).
Find angle above the negative \(x\) direction:
\(\theta=\arctan(0.400/2.30)\).
Write the commanded quadrant explicitly.}
{Use A/D for the vector root and C/D for division.
Use T, or interpolate the tangent table supplied on the reference sheet.
Signs define the quadrant; \(\theta\) is a small positive angle measured
from the negative \(x\) axis toward positive \(y\), not from positive \(x\).}
{Angle bands are GREEN within \(0.05^\circ\), AMBER through \(0.20^\circ\),
RED beyond. Wrong quadrant is always RED.
Magnitude and time still use ordinary relative-error bands.}
{AMBER: the predictor flags aim drift; charge 1 for an attitude reset.
RED: a wrong quadrant or angle predicts a missed entry corridor;
charge 2 for a complete manual-navigation rehearsal.}
{\(\Delta v=\sqrt{2.30^2+0.400^2}
=\sqrt{5.45}=\boxed{2.33\text{ m/s}}\).
\(t=2.3345235/0.155=\boxed{15.1\text{ s}}\).
Ratio \(=0.400/2.30=0.173913\).
T gives \(\boxed{9.87^\circ}\).
Table interpolation:
\(9.5+(0.173913-0.16734)/(0.17633-0.16734)(0.5)
=9.866^\circ\).
Command \(\boxed{\text{negative }x,\ \text{positive }y}\),
or about \(170.13^\circ\) counterclockwise from positive \(x\).
The historical maneuver's components are not being reconstructed.}

\Card{13-09}{Thread the entry corridor}
{Final historical correction lasted 22.4 s; game entry predictor is invented. [N1]}
{Swigert prepares Odyssey while Lovell checks the last correction.
The flight director supplies a locally linear response model.
It is deliberately simpler than a real Earth-entry trajectory.}
{Current modeled entry angle is \(-5.90^\circ\); target is
\(-6.50^\circ\). A positive corrective \(\Delta v\) changes angle by
\(-0.0400^\circ\) per m/s. Burn acceleration is 0.670 m/s\(^{2}\).
Find required \(\Delta v\), nominal burn duration, and predicted angle.
Find the minimum and maximum burn durations allowed by
\(-6.80^\circ\leq\gamma\leq-6.20^\circ\).}
{Use signed subtraction for \(\Delta\gamma\).
C/D: \(\Delta v=(\gamma_{\rm target}-\gamma_0)/k\) and
\(t=\Delta v/a\).
Repeat for both corridor edges.
Use \(\gamma(t)=\gamma_0+kat\) to check signs.}
{Command a positive correction with a negative angle response.
The nominal pad must target \(-6.50^\circ\), not merely any
duration in the corridor. Angle bands are \(0.05^\circ/0.20^\circ\).
Report the two duration bounds in ascending order.}
{AMBER: a small duration error produces a corridor warning; charge 1.
RED: a too-short burn predicts shallow skip-out; too long predicts
excessive entry loading. Charge 2 and correct before separation.}
{\(\Delta\gamma=-6.50-(-5.90)=-0.600^\circ\).
\(\Delta v=(-0.600)/(-0.0400)=\boxed{15.0\text{ m/s}}\).
\(t=15/0.670=\boxed{22.4\text{ s}}\);
\(\gamma=-5.90-0.0400(15)=\boxed{-6.50^\circ}\).
Shallow edge: \(\Delta v=7.50\) m/s,
\(t_{\min}=\boxed{11.2\text{ s}}\) (11.19403).
Steep edge: \(\Delta v=22.5\) m/s,
\(t_{\max}=\boxed{33.6\text{ s}}\) (33.58209).
Thus \(11.19403\leq t\leq33.58209\) s is the toy corridor.
The 22.4 s historical agreement is a chosen training calibration.}

\Card{13-10}{Wake Odyssey and come home}
{Historical LM jettison 141:30:02; splashdown 142:54:41. [N1]}
{Swigert powers Odyssey back up. Haise and Lovell check the final
energy plan before saying goodbye to Aquarius.
You must finish reactivation before jettisoning the lifeboat.}
{A separate, protected CM battery starts with 3.00 kWh.
Reactivation uses 0.800 kW for 45 min; entry 1.20 kW for 30 min;
recovery beacon 0.150 kW for 4.00 h.
Find each energy, total use, and energy remaining; retain 0.500 kWh.
Reactivation begins at 140:30:02: find completion GET and margin
until LM jettison. Find LM-jettison-to-splashdown interval.}
{Convert minutes to hours on C/D, then use \(E=Pt\).
Add energy with pencil; subtract from 3.00.
Add 45 min to reactivation GET and subtract the historical event times.
Beacon use extends after splashdown and is included in this pool.}
{Energy remaining must be at least 0.500 kWh and reactivation must
finish before 141:30:02. The CM battery was not charged for the
LM's earlier loads; do not merge those inventories.}
{AMBER: a sequencing error threatens the transfer; charge 1.
RED: premature jettison or insufficient entry energy fails the recovery
predictor; charge 2 and correct the sequence. If reserve remains
nonnegative and all gates are corrected, narrate parachutes and recovery
of all three crew. Otherwise end with rescue unproven.}
{Reactivation \(=0.800(45/60)=\boxed{0.600\text{ kWh}}\).
Entry \(=1.20(30/60)=\boxed{0.600\text{ kWh}}\).
Beacon \(=0.150(4)=\boxed{0.600\text{ kWh}}\).
Total \(=\boxed{1.80\text{ kWh}}\);
left \(=3.00-1.80=\boxed{1.20\text{ kWh}}\), GO.
Completion \(140{:}30{:}02+00{:}45{:}00
=\boxed{141{:}15{:}02}\).
Jettison margin \(=\boxed{00{:}15{:}00}\), GO.
Jettison to splashdown
\(142{:}54{:}41-141{:}30{:}02
=\boxed{01{:}24{:}39}\).
The power schedule is fictional; the separation and recovery anchors
come from the NASA catalog.}
